%% ============================================================
%%  Appendices
%% ============================================================

%% ─────────────────────────────────────────────────────────────
\chapter{Repository Map}
\label{app:repo}

\begin{lstlisting}[language=bash,
                   caption={Complete \medxplain{} repository structure.},
                   label={lst:repo}]
medxplain-simple/
|-- data_ingestion/              # Hospital-grade data pipeline
|   |-- dicom_pipeline.py        # DICOM loader -> DICOMStudy tensors
|   |-- phi_scrubber.py          # HIPAA/GDPR PHI removal + audit log
|   |-- report_parser.py         # Radiology report parser
|   `-- __init__.py
|-- training/
|   |-- train_vqa_two_phase.py   # Two-phase training script
|   |-- trainer.py               # Base trainer (AMP, DDP)
|   |-- losses.py                # NT-Xent, CE losses
|   `-- augmentations.py
|-- evaluation/
|   |-- evaluate_metrics_spec.py # VQA Acc, BLEU, BERTScore, Clin. F1
|   |-- clinical_validation.py   # Radiologist comparison + PDF reports
|   |-- benchmark.py
|   `-- visualization.py
|-- explainability/
|   |-- grad_cam.py              # Grad-CAM (pytorch-grad-cam)
|   |-- integrated_gradients.py  # Axiomatic attribution
|   `-- counterfactual.py
|-- medical_vqa_infrastructure/  # Architecture layer
|   |-- models/
|   |   |-- vision_encoder.py    # BiomedCLIP wrapper
|   |   |-- text_decoder.py      # BioGPT wrapper
|   |   |-- fusion.py            # Cross-attention fusion
|   |   `-- vqa_model.py         # Assembled VQA model
|   |-- training/
|   `-- config/
|-- vqa_app_deliverable/         # Clinical Gradio interface
|   |-- app_gradio.py            # Hospital-ready UI
|   |-- model_inference.py       # Grad-CAM + VQA inference
|   |-- advanced_features.py     # Clinical feature tabs
|   |-- config.py
|   `-- utils.py
|-- config/
|   `-- hospital_config.yaml     # Deployment configuration
|-- tests/
|   |-- test_dicom_pipeline.py   # DICOM pipeline unit tests
|   `-- test_phi_scrubber.py     # PHI scrubber unit tests
|-- docs/
|   |-- thesis_structure.md
|   `-- thesis/                  # This LaTeX thesis
|-- Readme.md
`-- requirements.txt
\end{lstlisting}

%% ─────────────────────────────────────────────────────────────
\chapter{Hospital Configuration Reference}
\label{app:config}

Listing~\ref{lst:config} reproduces the annotated
\code{hospital\_config.yaml} template.

\begin{lstlisting}[language=yaml,
                   caption={\code{config/hospital\_config.yaml} reference.},
                   label={lst:config}]
hospital:
  name: "[REDACTED]"          # Never commit real name
  data_root: "/data/dicom/"   # Override: MEDXPLAIN_DATA_ROOT

privacy:
  phi_scrubbing: true         # Must be true in production
  allow_cloud_upload: false   # Never upload to external servers
  audit_log_path: "./logs/phi_audit.log"
  audit_log_rotation_days: 90

model:
  vision_encoder: "biomedclip"  # biomedclip | resnet50 | densenet121
  language_model: "biogpt"      # biogpt | llama | gpt2
  fusion_type: "cross_attention"
  hidden_size: 768
  use_lora: true
  lora_rank: 8
  lora_alpha: 16
  temperature: 1.263

inference:
  batch_size: 1               # Process one study at a time
  device: "cuda:0"
  enable_explainability: true
  confidence_threshold: 0.50  # Below this -> flag for review

evaluation:
  discrepancy_report_dir: "./discrepancy_reports/"
  metrics: [accuracy, cohen_kappa, bleu, bertscore, clinical_f1]

demo:
  port: 7860
  server_name: "0.0.0.0"
  share: false                # Never true in hospital environment
\end{lstlisting}

%% ─────────────────────────────────────────────────────────────
\chapter{Radiologist Validation Form}
\label{app:form}

The following form is administered to radiologists \emph{before} the
model answer is revealed, to prevent anchoring bias.

\bigskip

\begin{mdframed}
\textbf{MedXplain Validation Study --- Radiologist Annotation Form}

\vspace{4mm}
\begin{tabular}{@{}ll}
Study ID (pseudonymised): & \underline{\hspace{6cm}} \\[2mm]
Modality:                 & \underline{\hspace{6cm}} \\[2mm]
Date:                     & \underline{\hspace{6cm}} \\
\end{tabular}

\vspace{4mm}
\textbf{Clinical Question:}\\
\underline{\hspace{12cm}}\\[1mm]
\underline{\hspace{12cm}}

\vspace{3mm}
\textbf{Radiologist Answer} (write before seeing model output):\\
\underline{\hspace{12cm}}\\[1mm]
\underline{\hspace{12cm}}

\vspace{2mm}
\textbf{Confidence} (circle one): \quad 1 \quad 2 \quad 3 \quad 4 \quad 5

\vspace{2mm}
\textbf{Comments:}\\
\underline{\hspace{12cm}}\\[1mm]
\underline{\hspace{12cm}}

\vspace{5mm}
\rule{\textwidth}{0.4pt}

\vspace{3mm}
\textit{[Model output revealed below --- complete above before reading]}

\vspace{3mm}
\textbf{Model Answer:} \underline{\hspace{8cm}}
\quad \textbf{Confidence:} \underline{\hspace{2cm}}\%

\vspace{3mm}
\textbf{Do you agree with the model?}
\quad $\square$~Yes \quad $\square$~No \quad $\square$~Partially

\textbf{Is the Grad-CAM explanation helpful?}
\quad $\square$~Yes \quad $\square$~No \quad $\square$~Not applicable

\textbf{Would you change your answer after seeing the model output?}
\quad $\square$~Yes \quad $\square$~No

\vspace{3mm}
\textbf{Notes on Grad-CAM:}\\
\underline{\hspace{12cm}}

\vspace{5mm}
Signature: \underline{\hspace{5cm}} \quad Date: \underline{\hspace{3cm}}
\end{mdframed}

%% ─────────────────────────────────────────────────────────────
\chapter{Evaluation Output Format}
\label{app:eval_output}

Listing~\ref{lst:eval_output} shows the formatted report generated by
\code{evaluation/evaluate\_metrics\_spec.py}.

\begin{lstlisting}[language=bash,
                   caption={Example output of the evaluation suite.},
                   label={lst:eval_output}]
========================================
MEDXPLAIN EVALUATION REPORT
========================================
VQA Accuracy:          68.2%  (Target: >65%)  v
BLEU-4:                0.28   (Target: >0.25) v
BERTScore F1:          0.78   (Target: >0.75) v
Clinical F1 (Macro):   0.63   (Target: >0.60) v
----------------------------------------
Per-Pathology F1:
  Pneumothorax:              0.71
  Pleural Effusion:          0.68
  Pneumonia:                 0.59
  Cardiomegaly:              0.54
  Atelectasis:               0.51
  Pulmonary Nodule:          0.48
  Pulmonary Oedema:          0.44
========================================
\end{lstlisting}
